\documentclass[10pt]{article}
\usepackage{amsmath,amssymb,amsthm,xcolor}
\usepackage[margin=0.85in]{geometry}
\definecolor{ResolutionGreen}{RGB}{0,100,55}
\newcommand{\VF}{\mathrm{VF}}\newcommand{\RV}{\mathrm{RV}}
\title{Section 3: proposed manuscript corrections}
\author{}\date{Unapplied review draft}
\begin{document}{\Large\bfseries Section 3: proposed manuscript corrections}\par\smallskip\textit{Unapplied review draft}\par\medskip
All text below is proposed. The working annotated manuscript is unchanged.
Numbering refers to the current manuscript; section 3 now occupies printed
pages 31--35 after the approved section 2 expansions.

\paragraph{Remark 3.15 (\texttt{GLZ:char}), R050.}
Replace only ``$\mathbb Z$-linear'' in the first sentence by
{\color{ResolutionGreen}``piecewise $\mathbb Z$-affine''}.
Keep the restriction on translation constants and the subsequent statement
about piecewise $\mathrm{GL}_k(\mathbb Z)$ transformations.

Independent check: on the nonzero RV group the map
\[
 t\longmapsto\begin{cases}t&\operatorname{vrv}(t)\geq0,\\
 t^{-1}&\operatorname{vrv}(t)<0\end{cases}
\]
is definable and contracts to $\gamma\mapsto|\gamma|$. This is not a
single affine function on a nontrivial ordered divisible group. If
$|\gamma|=n\gamma+c$ everywhere, substituting $0$ gives $c=0$;
substituting $a$ and $-a$ for $a>0$ then gives $a=na=-na$, a contradiction.
This counterexample concerns functions; it need not be a bijection.

\paragraph{After Notation 3.21 (\texttt{RVcat:tag}), R053.}
Immediately before ``The map $\mathbb L$ induces a surjective semiring
homomorphism'', insert:
\begin{quote}\color{ResolutionGreen}
Suppose that $\mathbb S$ is $(\VF,\Gamma)$-generated.
\end{quote}
The same hypothesis already appears in Theorem 3.22. The insertion makes
it apply to the preceding assertions about both lift homomorphisms.
The definition of the lift on individual objects needs no such hypothesis.

The referee's particular example uses a trivially valued ground field and
therefore does not meet the newly approved standing hypothesis. Here is
an independent replacement. Let $F$ be a small algebraically closed
\emph{nontrivially} valued subfield. Choose $\tau\in\RV$ with
$\operatorname{vrv}(\tau)\notin\Gamma(F)$, and use the real parameter
structure generated by $F$ and $\tau$. The VF part of its definable closure
is $F$, by Lemma 2.13 and algebraic closedness of $F$. Hence
$\tau^\sharp$ has no definable VF point. The singletons $\{\tau\}$ and
$\{1\}$ are isomorphic objects of $\RV[1]$, whereas their lifts cannot be
definably bijective: $1^\sharp$ contains the definable point $1$.
Definable bijections preserve the existence of a definable point, in both
directions. The ground value group remains nontrivial.

\paragraph{Notation 3.20, R051.}
Correct the spellings ``nonhomogenous'' and ``homogenous'' to
{\color{ResolutionGreen}``nonhomogeneous'' and ``homogeneous''}.

\paragraph{Theorem 3.22 (\texttt{main:prop:HK}), R054.}
In the measured characterization, replace $\int_+$ by
{\color{ResolutionGreen}$\int_+^\mu$}. Retain the measured congruence.

\paragraph{Formal verification boundary.}
The Lean implementation keeps the object lift separate from descent to
isomorphism classes, uses full original-base function graphs, and keeps
the measured relation $\mathbf P=0$ separate from the unmeasured relation
$\mathbf P-1=0$. The construction of all Grothendieck semirings and the two
full HK presentations is ongoing. A generic algebraic quotient lemma or
clean axiom list alone does not establish those paper statements.
\end{document}
